\documentclass[11pt]{article}
\usepackage{graphicx} % Required for inserting images
\usepackage{enumitem}
\usepackage{xcolor}
\usepackage{soul} % \hl{} highlights for review notes
\usepackage{hyperref}
\usepackage{subcaption}
\usepackage{comment}
\usepackage{amsmath}
\usepackage{amssymb}
\usepackage{bm}
\usepackage{float}
\usepackage{listings}
\usepackage{multirow}
\usepackage{booktabs}
\usepackage[capitalise]{cleveref}
\usepackage{tikz}
\usetikzlibrary{arrows.meta, patterns}

\usepackage[margin=1in]{geometry}
\emergencystretch=4em % Lets long \texttt names wrap instead of overflowing the margin

% Items still missing before submission are highlighted in red
\newcommand{\missing}[1]{\textcolor{red}{[MISSING: #1]}}

% Includes the figure if the file exists; otherwise draws a red placeholder box with the expected path
\newcommand{\fig}[2]{%
    \IfFileExists{#2}%
        {\includegraphics[width=#1]{#2}}%
        {\fcolorbox{red}{white}{\begin{minipage}[c][3cm][c]{\dimexpr #1-2\fboxsep-2\fboxrule\relax}\centering\color{red}\small MISSING FIGURE\\{\scriptsize\texttt{\detokenize{#2}}}\end{minipage}}}%
}

% Placeholder box for a figure that has not been taken yet: \missingfig{width}{description}
\newcommand{\missingfig}[2]{%
    \fcolorbox{red}{white}{\begin{minipage}[c][3cm][c]{\dimexpr #1-2\fboxsep-2\fboxrule\relax}\centering\color{red}\small MISSING: figure of #2\end{minipage}}%
}

\lstset{
    language=Matlab,
    basicstyle=\ttfamily\footnotesize,
    keywordstyle=\color{blue},
    commentstyle=\color{green!50!black},
    stringstyle=\color{purple},
    numbers=left,
    numberstyle=\tiny\color{gray},
    frame=single,
    breaklines=true,
    columns=fullflexible,
    keepspaces=true,
    literate={°}{{$^\circ$}}1, % listings cannot read UTF-8 characters directly
}

\title{EECE7324: Laboratory Assignment 5 (Module 6)}
\author{Andrew Benincasa and Pau Sabater}
\date{November 2026}

% Table rows generated by our scripts (\input would add an empty row after the last \hline)
\makeatletter\newcommand{\inputtable}[1]{\@@input #1 }\makeatother

\begin{document}

\maketitle

%=====================================================================
\section*{Part A: Filtering}
%=====================================================================
\subsection*{Transmitter}
Nothing changes in the transmitter: the noise filtering is done at the receiver. We reuse the Module 5 waveforms (\missing{list of \texttt{.WAVEFORM} files used: 16/32-PSK and 16/32-QAM, Gray}).

\subsection*{Receiver}
\begin{enumerate}
    \item \missing{position of the filtering block in \texttt{receiverSC.m} (section \texttt{Noise Filtering}, before \texttt{frameSynchronization()}) and why we filter before synchronization.}

    \item \missing{parameters of \texttt{bandpass()}~\cite{bandpass} and our filter: passband from $f_c$ and $B = (1+\beta)R_s$, and the new line 55.}

\begin{lstlisting}[caption={Our noise filter in \texttt{receiverSC.m}.}, label={lst:bandpass}]
yFiltered = NaN; % MISSING
\end{lstlisting}

    \item \missing{effect of the filter on the power spectrum, how we tuned the parameters, and the final filter details (passband, steepness, stopband attenuation).}

    \begin{figure}[H]
        \centering
        \fig{\linewidth}{figures/A/psd_filter.pdf}
        \caption{Power spectrum before and after the bandpass filter. \missing{capture}}
        \label{fig:psd-filter}
    \end{figure}

    \item \missing{setup and distance; BER of 16-PSK and 32-PSK over the air with and without filtering.}

    \begin{table}[H]
        \centering
        \small
        \begin{tabular}{c|cc|cc}
            \hline
            \multirow{2}{*}{\shortstack{Modulation\\order}} & \multicolumn{2}{c|}{BER} & \multicolumn{2}{c}{EVM (\%)} \\
             & No filter & Filter & No filter & Filter \\
            \hline
            \inputtable{tables/filter_psk.tex}
        \end{tabular}
        \caption{16-PSK and 32-PSK over the air, with and without the noise filter.}
        \label{tab:filter-psk}
    \end{table}

    \begin{figure}[H]
        \centering
        \fig{\linewidth}{figures/A/rx_psk_filter.pdf}
        \caption{Received 16-PSK and 32-PSK constellations without (left) and with (right) the noise filter.}
        \label{fig:rx-psk-filter}
    \end{figure}

    \item \missing{conclusion: is out-of-band noise and interference a major problem for our signals?}

    \item \missing{same study for 16-QAM and 32-QAM.}

    \begin{table}[H]
        \centering
        \small
        \begin{tabular}{c|cc|cc}
            \hline
            \multirow{2}{*}{\shortstack{Modulation\\order}} & \multicolumn{2}{c|}{BER} & \multicolumn{2}{c}{EVM (\%)} \\
             & No filter & Filter & No filter & Filter \\
            \hline
            \inputtable{tables/filter_qam.tex}
        \end{tabular}
        \caption{16-QAM and 32-QAM over the air, with and without the noise filter.}
        \label{tab:filter-qam}
    \end{table}

    \begin{figure}[H]
        \centering
        \fig{\linewidth}{figures/A/rx_qam_filter.pdf}
        \caption{Received 16-QAM and 32-QAM constellations without (left) and with (right) the noise filter.}
        \label{fig:rx-qam-filter}
    \end{figure}
\end{enumerate}

%=====================================================================
\section*{Part B: Equalization}
%=====================================================================
\subsection*{Transmitter}
Nothing changes in the transmitter: with our setup, pre-distortion does not make a difference.

\subsection*{Receiver}
\begin{enumerate}[resume]
    \item \missing{position of the equalization block in \texttt{receiverSC.m}, and on what it works (samples or symbols, pilot as training).}

    \item \missing{structure of \texttt{equalization.m}: inputs, LS estimation, ZF/MMSE branch, \texttt{filter()} with \texttt{e}.}

    We define two equalizers: Type 1, zero forcing (ZF), and Type 2, minimum mean square error (MMSE).

\subsubsection*{LS Channel Estimation}

    \item \missing{LS estimator $\mathbf{h} = (\mathbf{S}^H\mathbf{S})^{-1}\mathbf{S}^H\mathbf{r}$ with \texttt{'} and \texttt{\textbackslash}~\cite{ctranspose,mldivide}, line 17.}

\begin{lstlisting}[caption={LS channel estimation in \texttt{equalization.m}.}, label={lst:ls}]
h = NaN; % MISSING
\end{lstlisting}

    \item \missing{role of $L$ in the estimation; lengths of \texttt{h} and \texttt{H}.}

\subsubsection*{ZF Channel Equalization}

    \item \missing{ZF equalizer in the frequency domain, line 21; length of \texttt{e}.}

\begin{lstlisting}[caption={ZF equalizer in \texttt{equalization.m}.}, label={lst:zf}]
E = NaN; % MISSING
\end{lstlisting}

    \item \missing{constellation, BER and EVM of 16-PSK and 32-PSK with and without ZF for several $L$.}

    \begin{table}[H]
        \centering
        \small
        \begin{tabular}{c|c|cc}
            \hline
            $M$ & $L$ & BER & EVM (\%) \\
            \hline
            \inputtable{tables/zf_psk.tex}
        \end{tabular}
        \caption{16-PSK and 32-PSK without equalization ($L = $ --) and with ZF for several $L$.}
        \label{tab:zf-psk}
    \end{table}

    \begin{figure}[H]
        \centering
        \fig{\linewidth}{figures/B/rx_psk_zf.pdf}
        \caption{Received 16-PSK and 32-PSK constellations without equalization and with ZF for several $L$.}
        \label{fig:rx-psk-zf}
    \end{figure}

\subsubsection*{MMSE Channel Equalization}

    \item \missing{MMSE equalizer in the frequency domain, line 26, with \texttt{sigma2} and \texttt{Ps} (lines 23 and 24), \texttt{conj()} and \texttt{abs()}.}

\begin{lstlisting}[caption={MMSE equalizer in \texttt{equalization.m}.}, label={lst:mmse}]
E = NaN; % MISSING
\end{lstlisting}

    \item \missing{constellation, BER and EVM of 16-PSK and 32-PSK with and without MMSE for several $L$; comparison with ZF.}

    \begin{table}[H]
        \centering
        \small
        \begin{tabular}{c|c|cc}
            \hline
            $M$ & $L$ & BER & EVM (\%) \\
            \hline
            \inputtable{tables/mmse_psk.tex}
        \end{tabular}
        \caption{16-PSK and 32-PSK without equalization ($L = $ --) and with MMSE for several $L$.}
        \label{tab:mmse-psk}
    \end{table}

    \begin{figure}[H]
        \centering
        \fig{\linewidth}{figures/B/rx_psk_mmse.pdf}
        \caption{Received 16-PSK and 32-PSK constellations without equalization and with MMSE for several $L$.}
        \label{fig:rx-psk-mmse}
    \end{figure}

    \item \missing{same study (ZF and MMSE) for 16-QAM and 32-QAM.}

    \begin{table}[H]
        \centering
        \small
        \begin{tabular}{c|c|cc|cc}
            \hline
            \multirow{2}{*}{$M$} & \multirow{2}{*}{$L$} & \multicolumn{2}{c|}{ZF} & \multicolumn{2}{c}{MMSE} \\
             & & BER & EVM (\%) & BER & EVM (\%) \\
            \hline
            \inputtable{tables/eq_qam.tex}
        \end{tabular}
        \caption{16-QAM and 32-QAM without equalization ($L = $ --) and with ZF and MMSE for several $L$.}
        \label{tab:eq-qam}
    \end{table}

    \begin{figure}[H]
        \centering
        \fig{\linewidth}{figures/B/rx_qam_eq.pdf}
        \caption{Received 16-QAM and 32-QAM constellations without equalization and with ZF and MMSE.}
        \label{fig:rx-qam-eq}
    \end{figure}
\end{enumerate}

\begin{thebibliography}{9}
    \bibitem{bandpass} MathWorks, ``bandpass: Bandpass-filter signals,'' MATLAB Signal Processing Toolbox documentation (R2025b). [Online]. Available: \url{https://www.mathworks.com/help/signal/ref/bandpass.html}
    \bibitem{ctranspose} MathWorks, ``ctranspose, ': Complex conjugate transpose,'' MATLAB documentation (R2025b). [Online]. Available: \url{https://www.mathworks.com/help/matlab/ref/ctranspose.html}
    \bibitem{mldivide} MathWorks, ``mldivide, \textbackslash: Solve systems of linear equations $Ax = B$ for $x$,'' MATLAB documentation (R2025b). [Online]. Available: \url{https://www.mathworks.com/help/matlab/ref/mldivide.html}
\end{thebibliography}

\end{document}
